\documentclass[10pt]{article}
\usepackage[margin=0.85in]{geometry}
\usepackage{amsmath,amssymb,amsthm}
\usepackage{booktabs}
\usepackage{graphicx}
\usepackage{hyperref}

% ---------- notation ----------
\newcommand{\C}{\mathcal{C}}      % customers
\newcommand{\V}{\mathcal{V}}      % locations (depot + customers)
\newcommand{\K}{\mathcal{K}}      % vehicles
\renewcommand{\labelitemii}{$\circ$} 

\title{Capacity-Feasible Clustering for the CVRP on a Neutral-Atom Analog Device\\
\large DuQuantum $\times$ Pasqal Challenge --- Track A --- Model write-up}
\author{Team: Alex and Bob (Alex Wang, Christopher Thai)}
\date{Date: 10/25/2026}

\begin{document}
\maketitle

% =========================================================
\section{Problem formulation}
\label{sec:formulation}

We state the full Capacitated Vehicle Routing Problem (CVRP) first. Every later modeling choice (Sec.~\ref{sec:trackA} onward) is stated as a modification of this model.

\subsection{Sets}
\begin{itemize}
  \item $\C = \{1,\dots,N\}$: customers.
  \item $\V = \{0\} \cup \C$: all locations; $0$ denotes the depot.
  \item $\K = \{1,\dots,K\}$: vehicles (identical).
\end{itemize}

\subsection{Parameters}
\begin{itemize}
  \item $p_i = (x_i, y_i) \in [0,100]^2$, $i \in \V$: location of node $i$; the depot location $p_0$ is fixed at $(50, 50)$ before clustering.
   \begin{itemize}
        \item $x_i$: $x$-coordinate of node $i$.
        \item $y_i$: $y$-coordinate of node $i$.
    \end{itemize}
  \item $d_i \in \mathbb{Z}_{>0}$, $i \in \C$: demand of customer $i$, with $d_i \le Q$ (Assumption 5: no customer's demand exceeds the vehicle capacity).
  \item $Q \in \mathbb{Z}_{>0}$: capacity of every vehicle.
  \item $K \in \mathbb{Z}_{>0}$: fleet size (fixed; vehicles may be unused).
  \item $c_{ij} = \sqrt{(x_i-x_j)^2+(y_i-y_j)^2}$, $i,j \in \V$: travel cost; symmetric and satisfies the triangle inequality.
\end{itemize}

\subsection{Decision variables}
\begin{itemize}
  \item $a_{ik} \in \{0,1\}$, $i \in \C$, $k \in \K$: $a_{ik}=1$ iff vehicle $k$ serves customer $i$ (assignment / clustering).
  \item $e_{ijk} \in \{0,1\}$, $i,j \in \V$, $i \ne j$, $k \in \K$: $e_{ijk}=1$ iff vehicle $k$ travels directly from $i$ to $j$ (edge / routing).
\end{itemize}

\subsection{Objective}
Minimize the total distance driven by the fleet:
\begin{equation}
  \min_{a,\,e} \; \sum_{k \in \K} \sum_{i \in \V} \sum_{\substack{j \in \V \\ j \ne i}} c_{ij}\, e_{ijk}.
  \label{eq:obj}
\end{equation}

\subsection{Constraints}
\begin{align}
  \sum_{k \in \K} a_{ik} &= 1
    && \forall i \in \C
    && \text{(equality: each customer served exactly once)} \tag{C1}\label{c:assign}\\
  \sum_{i \in \C} d_i\, a_{ik} &\le Q
    && \forall k \in \K
    && \text{(inequality: capacity)} \tag{C2}\label{c:cap}\\
  \sum_{j \in \V,\, j \ne i} e_{ijk} &= a_{ik}
    && \forall i \in \C,\ k \in \K
    && \text{(equality: leave $i$ iff served by $k$)} \tag{C3}\label{c:out}\\
  \sum_{j \in \V,\, j \ne i} e_{jik} &= a_{ik}
    && \forall i \in \C,\ k \in \K
    && \text{(equality: enter $i$ iff served by $k$)} \tag{C4}\label{c:in}\\
  \sum_{j \in \C} e_{0jk} &\le 1
    && \forall k \in \K
    && \text{(inequality: at most one trip per vehicle)} \tag{C5}\label{c:depot}\\
  \sum_{\substack{i,j \in S \\ i \ne j}} e_{ijk} &\le |S| - 1
    && \forall S \subseteq \C,\ |S| \ge 2,\ k \in \K
    && \text{(inequality: subtour elimination)} \tag{C6}\label{c:subtour}\\
  e_{ijk},\ a_{ik} &\in \{0,1\}
    && && \text{(binary domains)} \tag{C7}\label{c:bin}
\end{align}

\paragraph{Remarks.}
(i) Flow balance at the depot, $\sum_{j} e_{0jk} = \sum_{j} e_{j0k}$, is implied by
\eqref{c:out}--\eqref{c:in}: every arc has one tail and one head, and customer in- and
out-degrees are equal, so the depot's are too.
(ii) \eqref{c:subtour} forbids any cycle among customers that avoids the depot; combined
with \eqref{c:out}--\eqref{c:depot}, each used vehicle drives exactly one closed tour
through the depot.
(iii) \eqref{c:subtour} has exponentially many members; it is used only for the
formulation and the classical reference, never on the device.
(iv) The two constraints that matter for the device have different types:
\eqref{c:assign} is an equality with a single target value, while \eqref{c:cap} is an
inequality satisfied by every load in $[0,Q]$. This asymmetry drives the encoding
choices in Sec.~\ref{sec:encoding}.

\subsection{Assumptions and where they are enforced}
\begin{center}
\begin{tabular}{@{}lll@{}}
\toprule
\# & Assumption (challenge brief, Sec.~2) & Enforced by \\
\midrule
1 & Depot location fixed before clustering & Parameter $p_0$ \\
2 & Symmetric costs, triangle inequality & Parameter $c_{ij}$ (Euclidean) \\
3 & Single vehicle, single visit, no split delivery & \eqref{c:assign}, \eqref{c:out}, \eqref{c:in} \\
4 & One trip per vehicle, no mid-route reload & \eqref{c:depot}, \eqref{c:subtour} \\
5 & $d_i \le Q$ for all $i$ & Instance generator (Sec.~\ref{sec:instances}) \\
6 & Fixed $K$; unused vehicles allowed at zero cost & Set $\K$; \eqref{c:depot} is $\le 1$, not $= 1$ \\
\bottomrule
\end{tabular}
\end{center}

\subsection{Capacity-feasible clustering}
\label{sec:trackA}
Under cluster-first/route-second, Track A decides only $a$; the routing variables $e$ and
\eqref{c:out}--\eqref{c:subtour} are solved afterwards by an exact classical TSP, used only
to evaluate the CVRP gap. This decomposition is a heuristic (brief, Pitfall~2.1).

\paragraph{Register fleet size.} Let $D=\sum_i d_i$ and
\[
  K_{\min}=\min\{K'\in\{1,\dots,K\} : \text{some } a \text{ with } K' \text{ vehicles satisfies \eqref{c:assign}--\eqref{c:cap}}\},
\]
the exact bin-packing number of the demands; $K_{\min}\ge\lceil D/Q\rceil$, and the inequality can be strict
(e.g.\ three demands of $2$ with $Q=3$ need three bins although $\lceil 6/3\rceil=2$). $K_{\min}$ is well defined
exactly when the instance is feasible, and an assignment using $K_{\min}$ vehicles is also a valid plan for
the full fleet, with vehicles $k>K_{\min}$ unused (Assumption~6). We compute $K_{\min}$ exactly per instance
(Sec.~\ref{sec:instances}) and model only $K'=K_{\min}$ vehicles, so the clustering model always has a feasible
solution and every feasible shot is a valid CVRP plan. The CVRP reference is computed with the full fleet $K$,
not $K_{\min}$.

\paragraph{Objective.} We minimize the within-vehicle pairwise distance:
\begin{equation}
  \min_{a}\; f(a)=\sum_{k=1}^{K'}\sum_{i<j} c_{ij}\,a_{ik}a_{jk}
  \quad \text{s.t.} \quad \eqref{c:assign},\ \eqref{c:cap},\ a_{ik} \in \{0,1\},\ k\le K'.
  \label{eq:cluster}
\end{equation}
Its weights are non-negative, so it maps onto positive Rydberg couplings with no linear term
(Sec.~\ref{sec:encoding}). For any non-negative pairwise objective, moving a customer to an empty vehicle never
increases its value, so some minimizer spreads customers over every available vehicle, whereas by the triangle inequality merging two tours (when capacity permits) never increases
route cost; sizing the register to $K_{\min}$ largely removes this mismatch and also minimizes the atom count.
On tier M (Sec.~\ref{sec:instances}), routing the optimal clustering of this model exactly gives a mean
decomposition gap of $3.3\%$ (median $2.3\%$) on the 14 of 60 instances with $K_{\min}=2$, against $11.7\%$ (median
$12.1\%$) when the register is built for the full fleet $K=3$; on the 46 instances with $K_{\min}=K=3$ the two coincide
(mean $3.3\%$, median $1.7\%$).
The model gap (Sec.~\ref{sec:results}) is measured against the optimum of \eqref{eq:cluster} with $K'=K_{\min}$.

% =========================================================
\section{Encoding and variable accounting}
\label{sec:encoding}

\subsection{What the register minimizes}
In the analog (Ising) mode the register evolves under
\begin{equation}
  \frac{H(t)}{\hbar}=\sum_m \frac{\Omega(t)}{2}\hat\sigma^x_m-\delta(t)\sum_m \hat n_m
  +\sum_{m<m'}U_{mm'}\,\hat n_m\hat n_{m'},\qquad U_{mm'}=\frac{C_6}{\lVert \mathbf r_m-\mathbf r_{m'}\rVert^6},
\end{equation}
where $\mathbf r_m$ is the position of atom $m$. With $\Omega=0$ at the end of the sweep every operator is
diagonal, so a measured bitstring $n\in\{0,1\}^M$ has the classical cost
\begin{equation}
  E(n)=-\delta\sum_m n_m+\sum_{m<m'}U_{mm'}\,n_m n_{m'} .
  \label{eq:register}
\end{equation}
This is a QUBO in which the linear coefficient $-\delta$ is set by the (global) laser detuning and is
\emph{identical for every atom}, while the couplings are set by the atom positions and are
\emph{strictly positive}, decay as $r^{-6}$, and cannot be made exactly zero; $C_6$ is fixed. A target
QUBO is therefore directly realizable only if all its couplings are $\ge 0$ and all its linear coefficients
are equal (up to positive scaling and a constant). With AnalogDevice, $U\le C_6/(5\,\mu\mathrm m)^6\approx 55.4$~rad/$\mu$s
and $|\delta|\le125.66$~rad/$\mu$s.

\subsection{Textbook formulation: binary slack}
Each capacity inequality becomes $L_k+\sum_{b} w_b s_{kb}=Q$ with $L_k=\sum_i d_i a_{ik}$, slack bits
$s_{kb}\in\{0,1\}$, weights $w_b=2^b$ and $m_Q=\lceil\log_2(Q+1)\rceil$ bits per vehicle, penalized by
$B\sum_k\big(L_k+\sum_b w_b s_{kb}-Q\big)^2$ alongside the assignment penalty. Using $b^2=b$ for binary variables, expanding gives
positive couplings $2Bd_id_j$, $2Bd_iw_b$ and $2Bw_bw_{b'}$ (each vehicle's atoms become fully connected)
and linear coefficients $B\,d_i(d_i-2Q)$ and $B\,w_b(w_b-2Q)$ that differ from atom to atom. It costs
$K m_Q$ extra atoms and is not realizable with a global detuning.

\subsection{Our formulation}
One atom per assignment variable $a_{ik}$, $k\le K'$, with
\begin{equation}
  H(a)=f(a)+\sum_i A_i\Big(1-\sum_{k}a_{ik}\Big)^2+\lambda_2\sum_{k}\big(L_k^2-2QL_k\big).
  \label{eq:qubo}
\end{equation}
\paragraph{Capacity penalty and the choice $\lambda_1=0$.}
The capacity term is a slack-free quadratic penalty. We start from the unbalanced penalization of
Monta\~nez-Barrera et al.~\cite{montanez2024unbalanced}, applied to the leftover room $h_k=Q-L_k$ of each vehicle, where $h_k\ge0$ means vehicle $k$ is within capacity,
\begin{equation}
  P(h_k)=-\lambda_1 h_k+\lambda_2 h_k^2,\qquad \lambda_1,\lambda_2\ge 0 .
\end{equation}
Substituting $h_k=Q-L_k$,
\begin{align*}
  P(h_k) &= -\lambda_1(Q-L_k)+\lambda_2(Q-L_k)^2\\
         &= -\lambda_1Q+\lambda_1L_k+\lambda_2Q^2-2\lambda_2QL_k+\lambda_2L_k^2\\
         &= \lambda_2L_k^2+(\lambda_1-2\lambda_2Q)L_k+(\lambda_2Q^2-\lambda_1Q),
\end{align*}
and summing over the $K'$ vehicles, where the constant appears once per vehicle,
\begin{equation}
  \sum_{k}P(h_k)=\lambda_2\sum_k L_k^2+(\lambda_1-2\lambda_2Q)\sum_k L_k+K'(\lambda_2Q^2-\lambda_1Q).
  \label{eq:unbalanced}
\end{equation}
Under \eqref{c:assign} the total load is fixed:
\begin{equation}
  \sum_k L_k=\sum_k\sum_i d_i a_{ik}=\sum_i d_i\sum_k a_{ik}=\sum_i d_i=D .
\end{equation}
Hence, for every assignment satisfying \eqref{c:assign}, \eqref{eq:unbalanced} equals
$\lambda_2\sum_k L_k^2+(\lambda_1-2\lambda_2Q)D+K'(\lambda_2Q^2-\lambda_1Q)$, and $\lambda_1$ enters only through
the constant $\lambda_1(D-K'Q)$. This constant is identical for every valid clustering, so $\lambda_1$ cannot
change their energy ordering. On states violating \eqref{c:assign}, $\lambda_1$ only adds $\lambda_1 d_i$ to the
linear coefficient of each $a_{ik}$, which the choice of $A_i$ below would absorb. For any $\lambda_1$, once the $A_i$ are large enough that every
ground state satisfies \eqref{c:assign}, $\lambda_1$ does not affect the ground state; we set $\lambda_1=0$, for which
Sec.~\ref{sec:weights} gives such a bound. This removes one hyperparameter without loss of generality.

By the same argument, the term $-2\lambda_2Q\sum_kL_k$ is also constant on valid clusterings. We keep it because,
up to a constant, it makes the capacity term the squared distance of each load from full capacity,
$\lambda_2\sum_k(L_k^2-2QL_k)=\lambda_2\sum_k(L_k-Q)^2-\lambda_2K'Q^2$, and the $A_i$ and $\delta$ below are derived with it.

On valid clusterings the penalty therefore acts only through $\lambda_2\sum_kL_k^2$. With $\sum_kL_k=D$ fixed,
\begin{equation}
  \sum_k L_k^2=\frac{D^2}{K'}+\sum_k\Big(L_k-\frac{D}{K'}\Big)^2 ,
\end{equation}
so it penalizes the spread of loads around their mean $D/K'$. Because $K'=K_{\min}\ge\lceil D/Q\rceil$,
the mean satisfies $D/K'\le Q$: perfectly balanced loads would be feasible. Integer demands and the
competing distance objective $f(a)$ prevent exact balance, so loads stay below $Q$ in most but not all
instances. Capacity is therefore soft, and we report the feasibility rate. In terms of atoms,
$L_k^2=\sum_i d_i^2a_{ik}+2\sum_{i<j}d_id_j\,a_{ik}a_{jk}$ contributes positive couplings $2\lambda_2 d_id_j$
within each vehicle and linear coefficients $\lambda_2 d_i(d_i-2Q)$.

Expanding with $a^2=a$, the linear coefficient of $a_{ik}$ is $-A_i+\lambda_2 d_i(d_i-2Q)$. Choosing
\begin{equation}
  A_i=A+\lambda_2\Big[d_i(d_i-2Q)-\min_j d_j(d_j-2Q)\Big]\ \ge A
\end{equation}
makes every linear coefficient equal to $-\delta$ with $\delta=A+\lambda_2\max_j d_j(2Q-d_j)$, so a global
detuning suffices and no DMM is needed. The resulting coefficients are:
\begin{center}
\begin{tabular}{@{}ll@{}}
\toprule
Term & Coefficient \\
\midrule
Linear, every $a_{ik}$ & $-\delta$ (identical) \\
$a_{ik}a_{il}$, same customer, $k\ne l$ & $2A_i$ \\
$a_{ik}a_{jk}$, same vehicle, $i\ne j$ & $c_{ij}+2\lambda_2 d_id_j$ \\
$a_{ik}a_{jl}$, $i\ne j$, $k\ne l$ & $0$ \\
Constant & $\sum_i A_i$ \\
\bottomrule
\end{tabular}
\end{center}
All couplings are non-negative. The target zeros cannot be realized exactly. On assignments satisfying
\eqref{c:assign}, however, only the difference between the same-vehicle and cross-vehicle couplings of each customer
pair affects the energy ordering, so cross-vehicle couplings need not vanish, only be equal (Sec.~\ref{sec:embedding}). Fixing a customer to a vehicle would break the vehicle-relabeling symmetry and save $K'$ atoms, but it would make the linear coefficients non-uniform again, so we do not do it.

\subsection{Penalty weights}
\label{sec:weights}
Let $\Gamma_i=\sum_{j\ne i}(c_{ij}+2\lambda_2 d_id_j)$.
 
\medskip\noindent\textbf{Claim.} \emph{If $\lambda_2\ge0$ and $\forall i$ we have}
\begin{equation}
  A_i>\max\big(\Gamma_i,\ \lambda_2 d_i(2Q-d_i)\big),
  \label{eq:Abound}
\end{equation}
\emph{then no ground state of \eqref{eq:qubo} violates \eqref{c:assign}.}
 
\begin{proof}
Let $a$ violate \eqref{c:assign} for some customer $i$, and let $s=\sum_k a_{ik}\ne1$. We exhibit a state
$a'$, differing from $a$ in the single variable $a_{ik}$, with $H(a')<H(a)$. Changing $a_{ik}$ alters only
three terms of \eqref{eq:qubo}: the assignment penalty of customer $i$, the pairs of $f$ that contain
$a_{ik}$, and the capacity term of vehicle $k$. Write $L_k$ for the load of vehicle $k$ in $a$.
 
\emph{Case $s=0$.} Choose any vehicle $k$ and set $a_{ik}=1$. The assignment penalty of customer $i$ drops
from $A_i$ to $0$. The objective $f$ increases by $\sum_{j\ne i}a_{jk}c_{ij}$. Since customer $i$ is not on
vehicle $k$, $L_k=\sum_{j\ne i}d_ja_{jk}$, and the capacity term changes by
\begin{align*}
  &\lambda_2\big[(L_k+d_i)^2-2Q(L_k+d_i)\big]-\lambda_2\big[L_k^2-2QL_k\big]\\
  &\quad=\lambda_2\big[L_k^2+2L_kd_i+d_i^2-2QL_k-2Qd_i-L_k^2+2QL_k\big]\\
  &\quad=2\lambda_2d_iL_k+\lambda_2d_i(d_i-2Q)\\
  &\quad=\sum_{j\ne i}a_{jk}\,2\lambda_2d_id_j+\lambda_2d_i(d_i-2Q).
\end{align*}
Because $a_{jk}\in\{0,1\}$ and $1\le d_i\le Q$, the total change satisfies
\[
  \Delta H=-A_i+\sum_{j\ne i}a_{jk}\big(c_{ij}+2\lambda_2d_id_j\big)+\lambda_2d_i(d_i-2Q)
  \le -A_i+\Gamma_i<0 .
\]
 
\emph{Case $s\ge2$.} Choose a vehicle $k$ with $a_{ik}=1$ and set $a_{ik}=0$. The assignment penalty of
customer $i$ changes by $A_i\big[(s-2)^2-(s-1)^2\big]=-A_i(2s-3)\le-A_i$. The objective $f$ changes by
$-\sum_{j\ne i}a_{jk}c_{ij}\le0$. The capacity term of vehicle $k$ changes by
\begin{align*}
  &\lambda_2\big[(L_k-d_i)^2-2Q(L_k-d_i)\big]-\lambda_2\big[L_k^2-2QL_k\big]\\
  &\quad=\lambda_2\big[L_k^2-2L_kd_i+d_i^2-2QL_k+2Qd_i-L_k^2+2QL_k\big]\\
  &\quad=\lambda_2\big[d_i^2-2L_kd_i+2Qd_i\big]\\
  &\quad=\lambda_2d_i\,(d_i-2L_k+2Q)\\
  &\quad\le\lambda_2d_i\,(d_i-2d_i+2Q)=\lambda_2d_i(2Q-d_i),
\end{align*}
where the inequality uses $L_k\ge d_i$ (customer $i$ is on vehicle $k$) and $\lambda_2 d_i\ge0$. Hence
$\Delta H\le-A_i+\lambda_2d_i(2Q-d_i)<0$.
 
In both cases $a$ is not a minimizer of $H$. Since $H$ is defined on the finite set $\{0,1\}^{NK'}$, a
minimizer exists, and by the above it satisfies \eqref{c:assign}. $\qedhere$
\end{proof}
 
Condition \eqref{eq:Abound} is sufficient but loose; a large $A_i$ compresses the objective couplings after scaling to the device.

\paragraph{Chosen weights.} We use $\lambda_2=10$ and $A=\Gamma_{\max}/2$, with $\Gamma_{\max}=\max_i\Gamma_i$. This
choice lies below the bound \eqref{eq:Abound} on all 100 tier-S instances tested, so it is validated empirically
rather than proven: exhaustive enumeration of all $2^{NK'}$ bitstrings gave a ground state satisfying
\eqref{c:assign} on every instance, with outcomes identical to the rigorous choice; with $A=\Gamma_{\max}/4$,
\eqref{c:assign} failed for $\lambda_2\ge12$. At $\lambda_2=10$ the ground state is the constrained optimum on 79
instances, feasible but suboptimal on 19, and over capacity on 2. No $\lambda_2\in\{6,8,10,12,15\}$ removes all
overloads: the quadratic penalty grows with load imbalance rather than with overload, so a clustering that
exceeds $Q$ by one unit costs only $O(\lambda_2 d_{\max})$ more than a balanced one, while larger $\lambda_2$
increases the number of suboptimal ground states. Relative to the rigorous bound, $A=\Gamma_{\max}/2$ reduces the
median coupling dynamic range from $18.5$ to $11.5$. Because $\lambda_2$ trades distance against demand, it is
retuned for each tier. (Script: \texttt{step1\_qubo\_check.py}; seeds $0$--$99$; results in \texttt{step1\_results.csv}.)
On both over-capacity instances ($\rho=0.93$), the constrained optimum is the first excited energy level of the
QUBO, so it is reachable by sampling and selected by the best-feasible-shot rule.

\paragraph{Tier M.} (Script: \texttt{step4a\_tierM\_weights.py}; seeds $0$--$59$.)
With $N=10$ the full QUBO has $2^{20}$ ($K'=2$) or $2^{30}$ ($K'=3$) states, so we enumerate the $K'^N$ one-hot
assignments, on which the assignment penalty vanishes, and check that the lowest one-hot state is a strict local
minimum of the full QUBO; it was on every instance and every $\lambda_2$ tested. For six $K'=2$ instances we enumerated
all $2^{20}$ states and confirmed that the ground state is one-hot with $A=\Gamma_{\max}/2$. The tier-S value
$\lambda_2=10$ does not transfer, and the right value depends on how tightly the $K'$ vehicles are packed. For $K'=3$
(46 of 60 instances, effective load $D/(K'Q)$ between $0.68$ and $0.93$), $\lambda_2=1$ makes the lowest one-hot state
optimal on 34, feasible but suboptimal on 8 and over capacity on 4; larger $\lambda_2$ removes overloads but quickly
adds suboptimal states (22 of 46 at $\lambda_2=3$). The $K'=2$ instances (14 of 60) are packed much more tightly
(effective load $0.75$--$1.00$, a direct consequence of choosing the smallest feasible fleet) and need $\lambda_2=35$ to
remove overloads, giving 8 optimal, 6 suboptimal and none over capacity. We therefore use $\lambda_2=1$ for $K'=3$ and
$\lambda_2=35$ for $K'=2$ on tier M. The stronger penalty raises the median coupling dynamic range from 54 to 237,
which makes the $K'=2$ tier-M registers harder to embed.

\paragraph{Device scaling.} Let $U_{\max}$ be the largest coupling coefficient of \eqref{eq:qubo}. Multiplying
\eqref{eq:qubo} by $\sigma=55.4/U_{\max}$~rad/$\mu$s places the strongest coupling at the 5~$\mu$m minimum spacing;
the final detuning is $\sigma\delta$, which must satisfy $\sigma\delta\le125.66$~rad/$\mu$s. Under \eqref{eq:Abound},
$U_{\max}=2\max_iA_i$, since every same-vehicle coupling satisfies $c_{ij}+2\lambda_2d_id_j\le\Gamma_i<A_i$; for
penalty weights below the bound, $U_{\max}$ must be computed explicitly. This mapping is used by the
coefficient-matching and weighted log fits; the ladder layout of Sec.~\ref{sec:embedding} sets its detuning from a validity
window instead.

\subsection{Variable accounting}
\begin{center}
\begin{tabular}{@{}lccccc@{}}
\toprule
Encoding & Count & S & M & L & Fits in 80? \\
\midrule
One-hot assignment, no constraints & $NK$ & 8 & 30 & 60 & yes \\
+ binary-encoded slack & $NK + K\lceil \log_2(Q+1) \rceil$ & 14 & 45 & 80 & L at the cap \\
Our formulation & $NK_{\min}$ & 8$^{\dagger}$ & 20--30 & 30--60 & yes \\
\bottomrule
\end{tabular}
\end{center}
{\small $^{\dagger}$Instances with $K_{\min}=1$ (4 atoms) are trivial and are excluded by the generator (Sec.~\ref{sec:instances}).}

The atom cap is necessary but not sufficient. The emulators bound what can be run in practice
(state-vector methods to roughly 20 atoms, emu-mps to roughly 30--60 at moderate entanglement), and the
couplings must be realizable in a plane: the one-hot structure requires, for each customer, a strongly
coupled $K'$-clique, and between vehicles, couplings that should vanish.
% TODO: per-instance K_min and atom counts once instances are generated.

\subsection{Decoding, feasibility and repair}
\label{sec:decoding}
A measured bitstring $n\in\{0,1\}^{NK'}$ is decoded as $a_{ik}=n_{iK'+k}$. A shot is \emph{feasible} if $a$ satisfies
\eqref{c:assign} and \eqref{c:cap}, and its cost is $f(a)$. Following the brief (Pitfall~9.1), we report two quantities
separately: \emph{raw}, the best feasible shot with no post-processing, and \emph{repaired}, the best shot after every
shot passes through a fixed, light repair. The repair (1) keeps, for a customer assigned to several vehicles, the one
with the smallest within-vehicle distance to its other customers; (2) places each unassigned customer, heaviest first,
on the vehicle with room that adds the least distance (on any vehicle if none has room); (3) while a vehicle is
overloaded, moves the single customer whose move to a vehicle with room adds the least distance, and declares the shot
infeasible if no such move exists. It costs $O(N^2K')$ per shot. The same decoding and repair are applied to a uniform
random sampler as a baseline. The best shot is the feasible shot with the lowest $f$; its clusters are routed by an
exact TSP to compute the CVRP gap.

% =========================================================
\section{Instances and reference values}
\label{sec:instances}
\paragraph{Feasibility certification (Task A2).} The brief's assumptions do not guarantee feasibility:
$D\le KQ$ is necessary but not sufficient, and deciding feasibility is a bin-packing problem. For each
seeded draw of demands we solve bin packing exactly with the $O(N2^N)$ subset dynamic program
($dp[S]$ = lexicographic minimum of bins used and current-bin load), which returns $K_{\min}$ and an
explicit assignment into $K_{\min}$ vehicles. A draw is accepted only if $2\le K_{\min}\le K$: $K_{\min}>K$ is
infeasible, and $K_{\min}=1$ is trivial. The certificate shipped with each instance is this assignment.
An independent validator checks every location lies in $[0,100]^2$, $1\le d_i\le Q$ (Assumption~5),
\eqref{c:assign} and \eqref{c:cap} on the certificate, that it uses exactly $K_{\min}$ non-empty vehicles with
$K_{\min}\le K$ (hence feasibility for the full fleet), and minimality of $K_{\min}$: either $K_{\min}=\lceil D/Q\rceil$,
which is a lower bound, or an exhaustive backtracking search confirms that $K_{\min}-1$ vehicles cannot pack
the demands. Both algorithms reproduce the brief's counterexample and agree on 2000 random packings.
\paragraph{Reference values.} Every reference is a proven optimum. On tier S we enumerate all assignments and route
each cluster by brute force. On tier M (seeds $0$--$59$) we compute the optimal depot tour of every customer subset with
the Held--Karp dynamic program, then the CVRP optimum with the full fleet $K$ by an exact set-partition dynamic program
over capacity-feasible subsets, in which each block is anchored on its lowest customer so that every partition is
generated once. On every instance this optimum matched an enumeration of all $K^N$ labelled assignments, and both the
CVRP optimum and the optimum of the clustering model were unique.

\paragraph{CPLEX cross-check (tier S).} As required by the brief (Sec.~9.3), we solved the 20 tier-S instances of
Table~\ref{tab:tierS} with CPLEX~22.2 (community edition, via \texttt{docplex}; time limit \texttt{cplex maxtime} $=60$~s
per solve) and compared with brute force. The clustering model was solved with the products $a_{ik}a_{jk}$ linearized by
$y_{ijk}\ge a_{ik}+a_{jk}-1$, $y_{ijk}\ge0$; the CVRP with the full fleet was solved in the three-index arc formulation
\eqref{c:assign}--\eqref{c:depot}, with Miller--Tucker--Zemlin load constraints in place of \eqref{c:subtour}. Every
solve ended with status ``integer optimal solution'', and CPLEX matched brute force on 20 of 20 instances for both
problems. (Script: \texttt{cplex\_check.py}; results in \texttt{cplex\_check.csv}.)
% TODO: generator details (geometries, rho bands, seeds), per-instance rho and d_max/Q.

% =========================================================
\section{Embedding}
\label{sec:embedding}
We use three embedders: the standard coefficient-matching fit, a weighted log fit, and a layout specialized to $K'=2$.

\subsection{Coefficient matching (general)}
This is the standard approach of Pulser's QUBO tutorial~\cite{pulserqubo}: choose atom positions so that the realized
couplings $U_{mm'}=C_6/\lVert\mathbf r_m-\mathbf r_{m'}\rVert^6$ reproduce the scaled off-diagonal QUBO coefficients
$\sigma Q_{mm'}$, by minimizing the Frobenius norm of the difference. We minimize
\[
  \sum_{m<m'}\big(U_{mm'}-\sigma Q_{mm'}\big)^2\Big/\sum_{m<m'}\big(\sigma Q_{mm'}\big)^2
\]
over all pairs, including those whose target is zero, using analytic gradients (L-BFGS-B), 20 random restarts, and
quadratic penalties for the 5~$\mu$m spacing and the 38~$\mu$m radius. Because the spacing penalty is soft, the
result is dilated uniformly until the minimum distance is exactly 5~$\mu$m; this scales every coupling by the same
factor, which is absorbed into $\sigma$ and the detuning. The method applies to any $K'$. Its fidelity is the
relative Frobenius error over all pairs,
\[
  \varepsilon_F=\lVert U/\sigma-Q_{\mathrm{off}}\rVert_F\big/\lVert Q_{\mathrm{off}}\rVert_F .
\]

\subsection{Weighted log fit (general)}
Coefficient matching weighs every coupling by its absolute error, so the many weak couplings barely influence the fit,
and it treats the target zeros as hard targets. The weighted log fit changes both. It minimizes
\[
  \frac{\sum_{\text{wanted}}\omega_{mm'}\big(\log U_{mm'}-\log\sigma Q_{mm'}\big)^2}{\sum_{\text{wanted}}\omega_{mm'}}
  +\frac{1}{n_0}\sum_{\text{target }0}\max\!\Big(0,\ \log U_{mm'}-\log\big(\tfrac12\min_{\text{wanted}}\sigma Q\big)\Big)^2 ,
\]
where the weights are $\omega=10$ for one-hot pairs (same customer), which enforce \eqref{c:assign}, and $\omega=1$
otherwise, and $n_0$ is the number of target-zero pairs. The log error treats strong and weak couplings alike, and the
hinge lets unwanted couplings be non-zero as long as they stay below half of the weakest wanted coupling, consistent
with the observation below that they need not vanish. Optimizer, restarts, penalties and dilation are as for coefficient
matching, and it applies to any $K'$. We report its fidelity with the same $\varepsilon_F$.

\subsection{Ladder layout ($K'=2$)}
\paragraph{Intuition.} Coefficient matching reproduced the wanted couplings well, but the cross-vehicle couplings,
which should be zero, remained large enough to reorder the low-energy states (Table~\ref{tab:static}). Two
observations led to the ladder. First, on any assignment satisfying \eqref{c:assign}, each customer has exactly one
excited atom, so for each customer pair exactly one of its four couplings is active. Shifting every cross-vehicle
coupling of a pair by the same constant therefore changes all valid energies equally: cross-vehicle couplings need to
be \emph{equal}, not zero. Second, equality can be forced by symmetry. If every customer's two atoms are mirror images
across one common line, the two cross-vehicle distances of any pair coincide, and so do the two same-vehicle
distances. Placing vehicle~1's atoms on one rail and vehicle~2's on a parallel rail, with a customer's two atoms at the
same position along the rails, achieves this.

\paragraph{Layout.} Vehicle $k$'s atoms lie on the rail $y=(k-\tfrac12)b$ with $b=5$~$\mu$m, and customer $i$'s two
atoms at position $\xi_i$ along both rails. For every customer pair, the same-vehicle, cross-vehicle and one-hot
couplings are, exactly,
\[
  u_{ij}=\frac{C_6}{|\xi_i-\xi_j|^6},\qquad v_{ij}=\frac{C_6}{\big(|\xi_i-\xi_j|^2+b^2\big)^3},\qquad
  U_{\mathrm{oh}}=\frac{C_6}{b^6}.
\]
For an assignment satisfying \eqref{c:assign}, with customer $i$ on vehicle $k_i$, the register energy
\eqref{eq:register} is
\[
  E=-\delta_rN+\sum_{i<j}v_{ij}+\sum_{i<j,\;k_i=k_j}\big(u_{ij}-v_{ij}\big),
\]
and on the same assignments the target \eqref{eq:qubo} equals $\sum_{i<j,\,k_i=k_j}w_{ij}$ plus a constant, with
$w_{ij}=c_{ij}+2\lambda_2d_id_j$. The ordering of valid clusterings is therefore reproduced exactly if
$u_{ij}-v_{ij}=t\,w_{ij}$ for a scale $t>0$. We fit the positions $\xi_i$ and $t$ by minimizing
$\sum_{i<j}\big(u_{ij}-v_{ij}-t\,w_{ij}\big)^2/(t\bar w)^2$ (Nelder--Mead, 40 random restarts), with penalties for the
5~$\mu$m spacing along a rail, the register radius, and the window condition of the claim below. Only couplings enter
the fit; no states are enumerated. Its fidelity is the relative valid-energy error over customer pairs,
$\varepsilon_L=\lVert(u-v)/t-w\rVert/\lVert w\rVert$.

\medskip\noindent\textbf{Claim.} \emph{Let $B_i=\sum_{j\ne i}u_{ij}$. If $\max_iB_i<\delta_r<U_{\mathrm{oh}}$, every
ground state of the register energy \eqref{eq:register} satisfies \eqref{c:assign}.}

\begin{proof}
If both atoms of a customer are excited, de-exciting one of them changes $E$ by $\delta_r-U_{\mathrm{oh}}$ minus its
(non-negative) couplings to the other excited atoms, which is at most $\delta_r-U_{\mathrm{oh}}<0$. Hence a minimizer has
no customer with two excited atoms. In such a minimizer, suppose customer $i$ has no excited atom. Every other
customer $j$ then has at most one excited atom, whose coupling to either atom of $i$ is $u_{ij}$ or $v_{ij}\le u_{ij}$.
Exciting one of $i$'s atoms therefore changes $E$ by at most $-\delta_r+B_i<0$, a contradiction. So every customer has
exactly one excited atom.
\end{proof}

\noindent We fit with $\max_iB_i\le0.8\,U_{\mathrm{oh}}$ and set $\delta_r$ to the midpoint of the window. Unlike
Sec.~\ref{sec:weights}, this guarantee holds for the realized register itself, including every residual coupling.
Capacity \eqref{c:cap} remains soft. The layout is defined only for $K'=2$: with three rails in a row, the outer
one-hot coupling is $2^6=64$ times weaker and the window closes, and with a triangle of atoms per customer the
cross-vehicle distances are no longer equal. Its one-dimensional coordinates also cannot represent arbitrary pairwise
weights.

\subsection{Static comparison and coverage}
Table~\ref{tab:static} compares the realized registers' ground states, found by enumerating all $2^{NK'}$ bitstrings
(evaluation only). For reference, the exact QUBO's ground state is the constrained optimum on 79\% of 100 tier-S
instances (Sec.~\ref{sec:weights}). The two fidelity metrics measure different things and are not comparable.
\begin{table}[h]
\centering
\resizebox{\textwidth}{!}{%
\begin{tabular}{@{}lccccccc@{}}
\toprule
Embedder & Fidelity (median) & Optimal & Feas., subopt. & Overload & C1 violated & Opt.\ in lowest 3 levels \\
\midrule
Coefficient matching & $\varepsilon_F=0.225$ & 7 & 7 & 5 & 1 & 10 \\
Weighted log fit & $\varepsilon_F=0.371$ & 10 & 10 & 0 & 0 & 11 \\
Ladder & $\varepsilon_L=0.55$ & 7 & 13 & 0 & 0 & 11 \\
\bottomrule
\end{tabular}}
\caption{Ground state of the realized register, tier S, 20 instances.}
\label{tab:static}
\end{table}
The coefficient-matching result depends on the local minimum the optimizer reaches. In a second environment (different
\texttt{numpy}/\texttt{scipy} versions), the same script and seeds gave 6 optimal, 11 feasible but suboptimal,
3 overloaded and no \eqref{c:assign} violations, while the ladder fit reproduced exactly. We report the run from the
environment pinned in \texttt{requirements.txt}.

The weighted log fit gives more correct static ground states than coefficient matching in both environments we ran
(10 vs.\ 7 and 8 vs.\ 6 optimal), but a lower probability of reaching the optimum in emulation (Table~\ref{tab:tierS}:
$0.26$ vs.\ $0.29$; $0.19$ vs.\ $0.25$ in the second environment). Static ground-state correctness therefore did not
predict dynamic performance, so we select embedders by shot metrics. By a rule fixed before the comparison (higher
$P(\text{reach})$, with raw feasibility as the tiebreaker when the two differ by less than $0.03$), coefficient
matching is our general embedder. The ladder covers all tier-S instances but only part of the larger tiers, so we
use the embedders as follows,
report results separately by $K'$ and by embedder, and do not change the instance generator to favor $K'=2$:
\begin{center}
\resizebox{\textwidth}{!}{%
\begin{tabular}{@{}ll@{}}
\toprule
Instances & Embedder \\
\midrule
Tier S (all $K'=2$) & Ladder (main), coefficient matching (general), weighted log fit (comparison) \\
Tier M, $K_{\min}=2$ (25 of 100 seeds) & Ladder, coefficient matching (control) \\
Tier M, $K_{\min}=3$ (75 of 100 seeds) & Coefficient matching \\
Tier L ($K_{\min}=2$ on 1 of 40 seeds) & Coefficient matching, if emulation allows \\
\bottomrule
\end{tabular}}
\end{center}

\subsection{Tier M: static check}
At 20 and 30 atoms the full spectrum cannot be enumerated, so we rank the $K'^N$ one-hot assignments by the energy of
the realized register, with the tier-M weights of Sec.~\ref{sec:weights}, and check whether the lowest one is a strict
local minimum of the full register energy, a necessary condition for it to be the ground state. Table~\ref{tab:tierMstatic}
compares the best feasible one-hot state with the model optimum and with the mean of a uniformly random feasible
clustering.
\begin{table}[h]
\centering
\resizebox{\textwidth}{!}{%
\begin{tabular}{@{}lccccccc@{}}
\toprule
Registers & Fidelity (median) & Lowest one-hot: opt.\ / subopt.\ / overload & Local min. & Gap of best feasible & Random feasible & Optimum rank \\
\midrule
$K'=2$, ladder (14) & $\varepsilon_L=0.84$ & 0 / 10 / 4 & 14/14 & 27.6\% & 29.8\% & top 11\% \\
$K'=2$, coefficient matching (14) & $\varepsilon_F=0.57$ & 0 / 4 / 10 & 0/14 & 31.0\% & 29.8\% & top 22\% \\
$K'=3$, coefficient matching (46) & $\varepsilon_F=0.75$ & 0 / 19 / 27 & 0/46 & 84.5\% & 81.8\% & top 5\% \\
\bottomrule
\end{tabular}}
\caption{Tier M, realized registers ranked over one-hot assignments (medians over instances).}
\label{tab:tierMstatic}
\end{table}

\noindent The realized tier-M registers retain almost none of the objective: their best feasible one-hot state is no
better than a random feasible clustering, the optimum never appears among the three lowest one-hot levels, and for
coefficient matching the lowest one-hot state is never even a local minimum, so the true ground state violates
\eqref{c:assign}. The ladder keeps \eqref{c:assign} by construction but cannot represent the weights: with
$\lambda_2=35$ the median coupling dynamic range is 237, and ten customers on a line cannot reproduce 45 pairwise
weights. Emulating these registers would therefore measure the embedding, not the device.

% =========================================================
\section{Pulse design}
\label{sec:pulse}
We use the baseline of the brief (Sec.~8.2), built with Pulser~\cite{pulser2022} on \texttt{AnalogDevice}: the
amplitude is a three-point interpolated ramp $0\to\Omega_{\mathrm{peak}}\to0$ with $\Omega_{\mathrm{peak}}=2\pi\times1.5$~MHz
$=9.42$~rad/$\mu$s, the detuning a three-point interpolated sweep $-\delta_f\to0\to\delta_f$, and the duration
$T=4~\mu$s. The final detuning $\delta_f$ is set by the embedding: the window midpoint $\delta_r$ for the ladder
($49.9$~rad/$\mu$s on tier S) and $\sigma\delta$ for coefficient matching. All values are within the device limits. The
pulse is not yet tuned. Tier S is emulated exactly with \texttt{pulser-simulation} (\texttt{QutipBackendV2}, state
vector, no noise); 1000 shots per instance are sampled from the exact final-state probabilities with fixed seeds, and
we also record the exact probability of the optimal clustering.
% TODO: pulse tuning; emu-sv / emu-mps settings (bond dimension) for tier M.

% =========================================================
\section{Results}
\label{sec:results}
\paragraph{Tier S, baseline pulse.} Table~\ref{tab:tierS} reports 20 tier-S instances (8 atoms, $K'=2$), 1000 shots
each, with $\varepsilon=0$ for reaching the reference. Both references are proven optima found by brute force: the
clustering model's optimum, and the CVRP optimum with the full fleet $K=2$ and exact routing.
\begin{table}[h]
\centering
\resizebox{\textwidth}{!}{%
\begin{tabular}{@{}lcccccc@{}}
\toprule
Method & Feasible (raw) & Best gap $0$ (raw) & $P(\text{reach})$ & Feasible (repaired) & Best gap $0$ (rep.) & Median CVRP gap \\
\midrule
Ladder & 96.1\% & 15/20 & 0.323 & 99.8\% & 19/20 & 0.0\% \\
Coefficient matching & 68.9\% & 16/20 & 0.292 & 94.7\% & 18/20 & 0.0\% \\
Weighted log fit & 70.3\% & 16/20 & 0.261 & 98.8\% & 20/20 & 0.0\% \\
Uniform random & 2.9\% & 20/20 & 0.008 & 90.2\% & 20/20 & 0.0\% \\
\bottomrule
\end{tabular}}
\caption{Tier S. Feasibility and $P(\text{reach})$ are averaged over instances; ``best gap $0$'' counts instances whose
best feasible shot is the clustering optimum. The CVRP gap uses the repaired best shot.}
\label{tab:tierS}
\end{table}

\noindent Tier S is too small to compare methods by best gap: each instance has only a handful of feasible
clusterings, so even the about 29 feasible shots among 1000 uniform random bitstrings contain the optimum on every
instance. The informative metrics are the feasibility rate and the probability of reaching the reference. There the
emulated register concentrates probability on good solutions: $P(\text{reach})$ is $0.32$ for the ladder and $0.29$ for
coefficient matching, against $0.008$ for random sampling, and 96\% and 69\% of shots are feasible against 3\%. The
ladder has the highest feasibility and $P(\text{reach})$, yet its best shot is optimal on fewer instances (15 vs.\ 16):
it concentrates on its realized ground state, which is the wrong clustering on 13 of 20 instances
(Table~\ref{tab:static}). Its exact probability of the optimum equals the sampled value, so the remaining loss comes
from the embedding, not from sampling or the sweep. The weighted log fit reaches the optimum with probability $0.26$,
below coefficient matching ($0.29$), despite having more correct static ground states (Table~\ref{tab:static}). Repair brings every method, including random sampling, close to
full feasibility, so on tier S repaired results reflect the repair rather than the device; we treat raw results as
primary. Tier M, where random sampling cannot find the optimum, is needed to compare methods by gap.

% =========================================================
\section{What did not work}
\label{sec:failures}
\begin{itemize}
  \item \textbf{Weak capacity penalty.} With $\lambda_2\le2$, 8 of 30 tier-S ground states exceeded capacity, all with
  load ratio $\rho\ge0.79$. No $\lambda_2\in\{6,8,10,12,15\}$ removed all overloads on 100 instances
  (Sec.~\ref{sec:weights}): the penalty measures load imbalance, not overload.
  \item \textbf{Too small a one-hot penalty.} With $A=\Gamma_{\max}/4$, the ground state violated \eqref{c:assign} for
  $\lambda_2\ge12$, so $A$ cannot be lowered much below $\Gamma_{\max}/2$.
  \item \textbf{Coefficient matching does not preserve the ground state.} Fitting every coupling, including the target
  zeros, by least squares (relative Frobenius error, median $0.225$) reproduced the wanted couplings to a median
  error of $2.4\%$, yet the realized ground state was the constrained optimum on only 7 of 20 instances, against 79\%
  for the exact QUBO, and was over capacity on 5 and violated \eqref{c:assign} on 1. The residual cross-vehicle couplings reached about half of the weakest wanted coupling,
  enough to reorder the low-energy states. Coefficient fidelity was therefore a poor predictor of ground-state
  correctness.
  \item \textbf{Forcing unwanted couplings small breaks the constraints.} A log-space fit that required every
  unwanted coupling to stay below 10\% of the weakest wanted one violated \eqref{c:assign} on all 20 instances: the
  geometry cannot satisfy the requirement, and the optimizer sacrificed the one-hot couplings. Weighting the
  one-hot couplings $10\times$ and relaxing the ceiling to 50\% restored \eqref{c:assign} but gave 10 of 20.
  \item \textbf{Static ground states as a selection criterion.} The weighted log fit produced more correct realized
  ground states than coefficient matching (10 vs.\ 7 of 20) but reached the optimum less often in emulation
  ($P(\text{reach})$ $0.26$ vs.\ $0.29$), and the same ordering held in a second environment ($0.19$ vs.\ $0.25$).
  Choosing it from static results alone would have been the wrong call, so embedders are now selected by shot metrics.
  \item \textbf{Direct embedding does not scale to tier M.} On 20- and 30-atom registers, both embedders lost
  essentially all objective information (Table~\ref{tab:tierMstatic}): the best feasible one-hot state was no better
  than a random feasible clustering. The stronger capacity penalty needed by tightly packed $K'=2$ instances makes the
  coupling dynamic range about 20 times larger than on tier S.
  \item \textbf{Limits of the ladder layout.} The two-rail layout (Sec.~\ref{sec:embedding}) guarantees
  \eqref{c:assign} and reproduced no overloads, but its one-dimensional customer coordinates cannot represent the
  target pairwise weights (median relative valid-energy error $0.55$), and it gave 7 of 20 optimal, within noise of
  the baselines. It is defined only for $K'=2$: all tier-S instances, 25 of 100 tier-M instances and 1 of 40 tier-L
  instances under our generator; $K'\ge3$ instances use coefficient matching.
\end{itemize}

% =========================================================
\section{Use of AI}
\label{sec:ai}
% TODO: where AI was used and for what. Include: the ladder-layout idea and its validity argument, and the
%       step1/step2/step2b/step3 scripts, were developed with AI assistance (Claude); state what each team member
%       re-derived or verified by hand.

% =========================================================
\begin{thebibliography}{9}
\bibitem{montanez2024unbalanced}
J.~A. Monta\~nez-Barrera, D.~Willsch, A.~Maldonado-Romo, and K.~Michielsen,
``Unbalanced penalization: a new approach to encode inequality constraints of combinatorial problems
for quantum optimization algorithms,''
\emph{Quantum Science and Technology} \textbf{9}, 025022 (2024).
arXiv:\href{https://arxiv.org/abs/2211.13914}{2211.13914}.
\bibitem{pulserqubo}
Pulser documentation, ``QAA to solve a QUBO problem,''
\url{https://pulser.readthedocs.io/en/stable/tutorials/qubo.html}.
\bibitem{pulser2022}
H.~Silv\'erio, S.~Grijalva, C.~Dalyac, L.~Leclerc, P.~J. Karalekas, N.~Shammah, M.~Beji, L.-P. Henry, and L.~Henriet,
``Pulser: An open-source package for the design of pulse sequences in programmable neutral-atom arrays,''
\emph{Quantum} \textbf{6}, 629 (2022).
\end{thebibliography}

\end{document}